\begin{problem}[第30届全国中学生物理竞赛复赛][刚体碰撞与转动]{2}
一长为 $2l$ 的轻质刚性细杆位于水平的光滑桌面上，杆的两端分别固定一质量为 $m$ 的小物块 $D$ 和一质量为 $\alpha m$（$\alpha$ 为常数）的小物块 $B$，杆可绕通过小物块 $B$ 所在端的竖直固定转轴无摩擦地转动。一质量为 $m$ 的小环 $C$ 套在细杆上（$C$ 与杆密接），可沿杆滑动，环 $C$ 与杆之间的摩擦可忽略。一轻质弹簧原长为 $l$，劲度系数为 $k$，两端分别与小环 $C$ 和物块 $B$ 相连。一质量为 $m$ 的小滑块 $A$ 在桌面上以垂直于杆的速度飞向物块 $D$，并与之发生完全弹性正碰，碰撞时间极短。碰撞时滑块 $C$ 恰好静止在距轴为 $r(r>l)$ 处。
\begin{subproblem}
若碰前滑块 $A$ 的速度为 $v_{0}$，求碰撞过程中轴受到的作用力的冲量；
\end{subproblem}
\begin{subproblem}
若碰后物块 $D$、$C$ 和杆刚好做匀速转动，求碰前滑块 $A$ 的速度 $v_{0}$ 应满足的条件。
\end{subproblem}
\end{problem}